\documentclass[aspectratio=169,10pt]{beamer}

\usetheme{SimplePlus}
\usepackage[T1]{fontenc}
\usepackage[utf8]{inputenc}
\usepackage{amsmath,amssymb}
\usepackage{graphicx}
\usepackage{booktabs}
\usepackage{hyperref}

\graphicspath{{../self-material/book/part-1/images/}}

\title{Macroeconomics PhD Intro\\Lecture 3: Dynamic Optimization Methods}
\subtitle{Chapter 4: Sequential and Recursive Methods}
\author{Swapnil Singh}
\institute{Lietuvos Bankas \and Kaunas University of Technology}
\date{Spring 2026}

\newcommand{\apxbutton}[1]{\hfill\hyperlink{#1}{\beamerbutton{Appendix}}}
\newcommand{\backbutton}[1]{\hfill\hyperlink{#1}{\beamerbutton{Back}}}

\begin{document}

\begin{frame}
  \titlepage
\end{frame}

\begin{frame}{Learning Goals}
\begin{itemize}
  \item Formulate dynamic optimization problems with states, controls, and feasibility sets.
  \item Solve finite-horizon models with Kuhn-Tucker conditions and Euler equations.
  \item Understand infinite-horizon subtleties: boundedness, no-Ponzi condition, and transversality.
  \item Derive and interpret Bellman equations and policy functions.
  \item Connect sequential and recursive methods in the neoclassical growth model.
  \item Build intuition for convergence and balanced growth under optimizing behavior.
\end{itemize}
\end{frame}

\begin{frame}{Road Map (Following Chapter 4)}
\begin{enumerate}
  \item Dynamic optimization setup and two benchmark models.
  \item Sequential methods with finite horizon.
  \item Sequential methods with infinite horizon.
  \item Recursive methods and the Bellman equation.
  \item Dynamics and convergence in the optimizing neoclassical growth model.
\end{enumerate}

\begin{block}{Connection to Lecture 2}
In Solow, saving was exogenous. In this lecture, saving is chosen optimally from primitives.
\end{block}
\end{frame}

%% ============================================================
\section{Motivation}
%% ============================================================

\begin{frame}{Why Dynamic Optimization?}
Cass-Koopmans adds micro-foundations to growth: saving/investment become outcomes of rational choice.

\medskip
\begin{itemize}
  \item Solow: $i_t=s y_t$ imposed.
  \item Optimizing growth: $i_t$ implied by Euler equation.
\end{itemize}

\begin{block}{What Changes?}
\begin{itemize}
  \item Positive analysis: richer transition dynamics and policy propagation.
  \item Normative analysis: welfare comparisons become internally coherent.
\end{itemize}
\end{block}
\end{frame}

\begin{frame}{Sequential vs Recursive: Conceptual Difference}
\begin{columns}[T]
\column{0.49\textwidth}
\textbf{Sequential}
\begin{itemize}
  \item Unknown: entire sequence.
  \item Method: Kuhn-Tucker/FOCs over time.
  \item Output: path $\{x_t,y_t\}$.
\end{itemize}

\column{0.49\textwidth}
\textbf{Recursive}
\begin{itemize}
  \item Unknown: value function $V(\cdot)$.
  \item Method: functional equation.
  \item Output: policy rule $x'=g(x)$.
\end{itemize}
\end{columns}

\medskip
Both represent the same economics under standard assumptions.
\end{frame}

%% ============================================================
\section{4.1 Problem Setup}
%% ============================================================

\begin{frame}{Generic Dynamic Optimization Problem (P1)}
\hypertarget{m-a1}{}
\[
\max_{\{y_t,x_{t+1}\}_{t=0}^{T}} \sum_{t=0}^{T} \beta^t \hat{\mathcal{F}}(y_t)
\]
subject to
\[
x_{t+1}=h(x_t,y_t), \qquad x_{t+1}\in\Gamma(x_t)
\]

\begin{itemize}
  \item State: $x_t$ (predetermined stock variable).
  \item Control: $y_t$ (choice in period $t$).
  \item Initial condition: $x_0$ given.
\end{itemize}
\apxbutton{a1}
\end{frame}

\begin{frame}{Discounting and Intertemporal Trade-offs}
\begin{itemize}
  \item $\beta\in(0,1)$: stationary discount factor.
  \item Relative welfare weight between dates $t$ and $t+k$: $\beta^k$.
  \item Dynamic trade-off: reducing current payoff can raise future payoff through state dynamics.
\end{itemize}

\begin{block}{Interpretation}
The model separates \textbf{technology/constraints} $\{h,\Gamma\}$ from \textbf{preferences/objective} $\{\hat{\mathcal{F}},\beta\}$.
\end{block}
\end{frame}

\begin{frame}{Existence Intuition}
Sufficient conditions used in Chapter 4:
\begin{itemize}
  \item Continuity of objective and transition mappings.
  \item Feasible set nonempty, closed, bounded.
  \item Concavity/convexity structure.
\end{itemize}

\medskip
\textbf{Why it matters:}
\begin{itemize}
  \item A maximum exists (not just a supremum).
  \item FOCs characterize a unique global solution when problem is convex.
\end{itemize}
\end{frame}

\begin{frame}{State-Control Timeline}
Given $x_t$, choosing $y_t$ determines $x_{t+1}=h(x_t,y_t)$.

\medskip
\begin{itemize}
  \item At $t=0$: $x_0$ exogenous, choose $y_0$, get $x_1$.
  \item At $t=1$: $x_1$ now predetermined, choose $y_1$, get $x_2$.
  \item Continue recursively.
\end{itemize}

\begin{block}{Key Distinction}
Controls are current choices. States are consequences of past choices.
\end{block}
\end{frame}

%% ============================================================
\section{4.1.1 Consumption-Saving Model}
%% ============================================================

\begin{frame}{Consumption-Saving Model: Primitives}
\begin{itemize}
  \item Preferences: $U=\sum_{t=0}^{T}\beta^t u(c_t)$.
  \item $u'(c)>0$, $u''(c)<0$, and $\lim_{c\to 0}u'(c)=\infty$.
  \item Endowment income: wage sequence $\{w_t\}$.
  \item Assets $a_t$ with return $r_t$.
\end{itemize}

Budget constraint:
\[
c_t=w_t+(1+r_t)a_t-a_{t+1}
\]

Borrowing limit and terminal condition:
\[
a_{t+1}\ge \underline{a},\qquad a_{T+1}\ge 0
\]
\end{frame}

\begin{frame}{Economic Interpretation of Asset Dynamics}
\begin{itemize}
  \item If $a_{t+1}>a_t(1+r_t)+w_t$, consumption is negative (infeasible).
  \item $a_{t+1}<0$ means borrowing from the future.
  \item Terminal restriction avoids ``borrow and default'' in the last period.
\end{itemize}

\begin{block}{Trade-off}
One unit less consumption today buys $1+r_{t+1}$ more resources tomorrow.
\end{block}
\end{frame}

\begin{frame}{Consumption-Saving Problem (P2)}
\[
\max_{\{c_t,a_{t+1}\}_{t=0}^{T}}\sum_{t=0}^{T}\beta^t u(c_t)
\]
subject to
\[
c_t=w_t+(1+r_t)a_t-a_{t+1}
\]
\[
c_t\ge 0,\quad a_{t+1}\ge \underline{a},\quad a_{T+1}\ge 0
\]

Mapping to (P1): state $x_t=a_t$, control $y_t=c_t$.
\end{frame}

%% ============================================================
\section{4.1.2 Neoclassical Growth Model}
%% ============================================================

\begin{frame}{Neoclassical Growth Model: Setup}
Production economy (not pure endowment):
\[
y_t=f(k_t),\qquad k_{t+1}=(1-\delta)k_t+i_t
\]
Resource constraint:
\[
c_t+i_t\le y_t
\]
Equivalent form:
\[
c_t+k_{t+1}\le f(k_t)+(1-\delta)k_t
\]

\medskip
Same preferences as consumption-saving model, but return is endogenous through $f'(k)$.
\end{frame}

\begin{frame}{Planner Problem in the NGM (4.3)}
\[
\max_{\{c_t,k_{t+1}\}_{t=0}^{T}}\sum_{t=0}^{T}\beta^t u(c_t)
\]
subject to
\[
c_t+k_{t+1}\le f(k_t)+(1-\delta)k_t,
\qquad c_t,k_{t+1}\ge 0
\]

\begin{block}{Interpretation}
Competitive equilibrium with complete markets and perfect competition will decentralize this planner allocation (Chapter 5-6 logic).
\end{block}
\end{frame}

%% ============================================================
\section{4.2 Sequential Methods: Finite Horizon}
%% ============================================================

\begin{frame}{Finite Horizon: Strategy}
When $T<\infty$, use Kuhn-Tucker conditions.

\medskip
\textbf{Why this is tractable:}
\begin{itemize}
  \item Finite number of unknowns.
  \item Boundary condition at the terminal date.
  \item Concavity gives sufficiency of FOCs.
\end{itemize}
\end{frame}

\begin{frame}{Two-Period Consumption-Saving Example}
\hypertarget{m-a2}{}
With $a_0=0$, $\underline{a}=0$, and constant $r$:
\[
c_0=w_0-a_1,\qquad c_1=w_1+(1+r)a_1
\]
Lifetime budget:
\[
c_0+\frac{c_1}{1+r}=w_0+\frac{w_1}{1+r}
\]
Borrowing/saving restriction: $a_1\ge 0 \Leftrightarrow w_0-c_0\ge 0$.
\apxbutton{a2}
\end{frame}

\begin{frame}{Two-Period Lagrangian and FOCs}
Lagrangian:
\[
\mathcal{L}=u(c_0)+\beta u(c_1)+\mu\left(w_0+\frac{w_1}{1+r}-c_0-\frac{c_1}{1+r}\right)+\lambda(w_0-c_0)
\]
FOCs:
\[
u'(c_0)-\mu-\lambda=0,
\qquad
\beta u'(c_1)-\frac{\mu}{1+r}=0
\]
Complementary slackness:
\[
\lambda(w_0-c_0)=0,\quad w_0-c_0\ge 0,\quad \lambda\ge 0
\]
\end{frame}

\begin{frame}{Euler Equation and Cases}
If interior ($\lambda=0$):
\[
u'(c_0)=\beta(1+r)u'(c_1)
\]

\begin{itemize}
  \item $\beta(1+r)=1$: strongest smoothing benchmark.
  \item High current income and borrowing constraint can generate corner solutions.
  \item If borrowing bound binds, Euler inequality replaces equality.
\end{itemize}
\end{frame}

\begin{frame}{Generic T-Period Result (Result 1)}
Rewrite objective as
\[
\max_{\{x_{t+1}\}}\sum_{t=0}^{T}\beta^t\mathcal{F}(x_t,x_{t+1}),\qquad x_{t+1}\in[\underline{x},\gamma(x_t)]
\]
Interior Euler equation for $t<T$:
\[
\mathcal{F}_2(x_t^*,x_{t+1}^*)+\beta\mathcal{F}_1(x_{t+1}^*,x_{t+2}^*)=0
\]
Terminal condition:
\[
x_{T+1}^*=\underline{x}
\]

\medskip
This is a second-order difference system in the state sequence.
\end{frame}

\begin{frame}{Finite-Horizon NGM in Sequence Form (P3)}
\hypertarget{m-a3}{}
\[
\max_{\{k_{t+1}\}_{t=0}^{T}}\sum_{t=0}^{T}\beta^t u\left(f(k_t)+(1-\delta)k_t-k_{t+1}\right)
\]
subject to
\[
k_{t+1}\ge 0
\]

FOC for $t<T$:
\[
-u'(c_t)+\beta u'(c_{t+1})\left[f'(k_{t+1})+1-\delta\right]+\mu_t=0
\]
\apxbutton{a3}
\end{frame}

\begin{frame}{Terminal Condition in Finite-Horizon NGM}
Final-period FOC:
\[
-u'(c_T)+\mu_T=0
\]
Because $u'(c_T)>0$, we get $\mu_T>0$, hence by complementary slackness:
\[
k_{T+1}=0
\]

\begin{block}{Intuition}
Capital after the terminal date has no continuation value, so it is never optimal to leave it unconsumed.
\end{block}
\end{frame}

\begin{frame}{Euler Equation for NGM (Interior Dates)}
For $t<T$, Inada conditions imply interior $k_{t+1}>0$ and $\mu_t=0$:
\[
 u'(c_t)=\beta u'(c_{t+1})\left[f'(k_{t+1})+1-\delta\right]
\]
with
\[
c_t=f(k_t)+(1-\delta)k_t-k_{t+1}
\]

Unknowns pinned down by: Euler equations + $k_0$ + terminal condition $k_{T+1}=0$.
\end{frame}

\begin{frame}{Example 4.1: Closed-Form Finite-Horizon Solution}
Assume:
\[
u(c)=\log c,\qquad f(k)=Ak^{\alpha},\qquad \delta=1
\]
Then
\[
k_{t+1}=\alpha\beta\frac{1-(\alpha\beta)^{T-t}}{1-(\alpha\beta)^{T-t+1}}Ak_t^{\alpha}
\]
and
\[
c_t=\frac{1-\alpha\beta}{1-(\alpha\beta)^{T-t+1}}Ak_t^{\alpha}
\]

\textbf{Message:} optimal saving rate is time-varying because horizon is finite.
\end{frame}

\begin{frame}{Finite-Horizon Intuition: Time-to-Go Effect}
\begin{itemize}
  \item Far from terminal date: strong incentive to invest for future periods.
  \item Near terminal date: investment motive weakens.
  \item Last period: invest nothing beyond horizon.
\end{itemize}

\begin{block}{Contrast with Infinite Horizon}
In stationary infinite-horizon models, rules are time-invariant and depend on the state, not calendar time.
\end{block}
\end{frame}

%% ============================================================
\section{CRRA and EIS}
%% ============================================================

\begin{frame}{CRRA Utility Class}
\[
u(c)=\frac{c^{1-\sigma}-1}{1-\sigma},\qquad \sigma\ge 0,\;\sigma\ne 1
\]
Special cases:
\begin{itemize}
  \item $\sigma\to 1$: log utility.
  \item $\sigma>0$: strictly concave preferences.
  \item $\sigma=0$: linear utility.
\end{itemize}

CRRA is central because it delivers constant EIS and is compatible with balanced growth.
\end{frame}

\begin{frame}{Example 4.2: EIS Under CRRA}
\hypertarget{m-a4}{}
With linear returns $R_{t,t+s}$, repeated Euler substitution gives
\[
\frac{c_{t+s}}{c_t}=\left(\beta^sR_{t,t+s}\right)^{1/\sigma}
\]
Hence
\[
EIS=\frac{\partial\log(c_{t+s}/c_t)}{\partial\log R_{t,t+s}}=\frac{1}{\sigma}
\]

\begin{itemize}
  \item High $\sigma$: lower willingness to tilt consumption intertemporally.
  \item Low $\sigma$: stronger response to return changes.
\end{itemize}
\apxbutton{a4}
\end{frame}

%% ============================================================
\section{4.3 Sequential Methods: Infinite Horizon}
%% ============================================================

\begin{frame}{Why Infinite Horizon?}
\begin{itemize}
  \item Stationarity: problem at $t$ has same structure as at $t+1$.
  \item Useful approximation for long-run macro questions.
  \item Dynastic interpretation: current generation values descendants.
\end{itemize}

\medskip
But infinite horizon requires extra mathematical discipline.
\end{frame}

\begin{frame}{Infinite-Horizon Pitfall 1: Unbounded Utility}
Even with $\beta<1$, utility may be ill-defined if feasible consumption grows too fast.

For CRRA with growth $c_t=c_0(1+\gamma)^t$, boundedness needs:
\[
\beta(1+\gamma)^{1-\sigma}<1
\]

\begin{block}{Takeaway}
Discounting alone is not enough; curvature and growth jointly determine boundedness.
\end{block}
\end{frame}

\begin{frame}{Infinite-Horizon Pitfall 2: Too-Large Feasible Set}
Without borrowing limits, households can increase $c_0$ by rolling debt forever.

\medskip
Need no-Ponzi condition (nPg):
\[
\lim_{t\to\infty}\frac{a_{t+1}}{(1+r)^t}\ge 0
\]

\begin{itemize}
  \item Excludes exploding debt in present value.
  \item Makes optimization problem well posed.
\end{itemize}
\end{frame}

\begin{frame}{nPg and Lifetime Budget Constraint}
\hypertarget{m-a5}{}
From period-by-period budget constraints:
\[
c_t+a_{t+1}=(1+r)a_t
\]
Iterating forward and using nPg gives
\[
\sum_{t=0}^{\infty}\frac{c_t}{(1+r)^t}\le a_0(1+r)
\]

At optimum with strictly increasing utility, equality holds.
\apxbutton{a5}
\end{frame}

\begin{frame}{TVC: Distinct from nPg}
Transversality condition:
\[
\lim_{t\to\infty}\frac{a_{t+1}}{(1+r)^t}\le 0
\]

\begin{itemize}
  \item nPg: external feasibility restriction (no endless unpaid debt).
  \item TVC: internal optimality restriction (no over-accumulating assets that could be consumed).
\end{itemize}

Together with strict monotonicity of utility, they imply a tight lifetime budget equality.
\end{frame}

\begin{frame}{Example 4.3: Infinite-Horizon Log Consumption-Saving}
Problem:
\[
\max\sum_{t=0}^{\infty}\beta^t\log c_t
\quad\text{s.t.}\quad
c_t+a_{t+1}=(1+r)a_t
\]
Consolidated budget and FOCs imply:
\[
c_t=[\beta(1+r)]^t c_0,
\qquad
c_0=a_0(1-\beta)(1+r)
\]

This is a full closed-form consumption path.
\end{frame}

\begin{frame}{Proposition 4.4 (Sufficiency with TVC)}
\hypertarget{m-a6}{}
For
\[
\max_{\{x_{t+1}\}}\sum_{t=0}^{\infty}\beta^t\mathcal{F}(x_t,x_{t+1}),\quad x_{t+1}\in\Gamma(x_t)
\]
if
\begin{itemize}
  \item Euler equations hold at all $t$,
  \item TVC holds,
  \item $\mathcal{F}$ jointly concave and increasing in first argument,
  \item $\Gamma$ convex,
\end{itemize}
then the candidate sequence is optimal.
\apxbutton{a6}
\end{frame}

\begin{frame}{TVC in the Infinite-Horizon NGM}
In NGM:
\[
\mathcal{F}_1 \to u'(c_t)\left[f'(k_t)+1-\delta\right]
\]
TVC becomes
\[
\lim_{t\to\infty}\beta^t u'(c_t)\left[f'(k_t)+1-\delta\right]k_t=0
\]

\begin{block}{Economic Meaning}
The discounted shadow value of remaining capital cannot stay positive in the limit.
\end{block}
\end{frame}

\begin{frame}{Infinite-Horizon NGM: Steady State Condition}
At a stationary steady state ($c_{t+1}=c_t$, $k_{t+1}=k_t=\bar{k}$), Euler implies
\[
1=\beta\left[f'(\bar{k})+1-\delta\right]
\]

Then
\[
\bar{c}=f(\bar{k})-\delta\bar{k}
\]

This is the optimizing counterpart to Solow steady-state logic.
\end{frame}

\begin{frame}{Examples 4.5-4.6: Link to Solow}
For log utility, Cobb-Douglas, full depreciation:
\[
k_{t+1}=\alpha\beta A k_t^{\alpha}
\]
Hence saving is a constant share:
\[
s=\alpha\beta
\]

\begin{block}{Important Bridge}
Under these assumptions, optimizing growth generates Solow-style law of motion with endogenous $s$.
\end{block}
\end{frame}

\begin{frame}{Balanced Growth Version of NGM}
With $Y_t=F(K_t,A_tL_t)$, $A_t$ grows at $\gamma$, $L_t$ at $n$, and utilitarian welfare:
\[
\max \sum_{t=0}^{\infty}\beta^t L_t u(c_t)
\]

Normalize by efficiency units:
\[
\tilde{k}_t=\frac{K_t}{A_tL_t},\qquad \tilde{c}_t=\frac{C_t}{A_tL_t}
\]

Balanced growth path corresponds to steady values of normalized variables.
\end{frame}

\begin{frame}{Why CRRA is Central for Balanced Growth}
Chapter 4 emphasizes a sharp characterization:
\begin{itemize}
  \item Exact balanced growth with constant interest in these models requires power utility (CRRA/log limit).
  \item Otherwise, saving behavior drifts systematically with growth in scale.
\end{itemize}

\begin{block}{Practical Implication}
Most quantitative macro models use CRRA not only for convenience, but for consistency with long-run growth facts.
\end{block}
\end{frame}

%% ============================================================
\section{4.4 Recursive Methods}
%% ============================================================

\begin{frame}{Dynamic Programming Idea}
\hypertarget{m-a7}{}
Instead of choosing full sequences directly, define the value of starting from state $x$:
\[
V(x)=\max_{\{x_{t+1}\}}\sum_{t=0}^{\infty}\beta^t\mathcal{F}(x_t,x_{t+1})
\]

Stationarity allows decomposition into today + continuation value.
\apxbutton{a7}
\end{frame}

\begin{frame}{Bellman Equation (4.22)}
\[
V(x)=\max_{x'\in\Gamma(x)}\left\{\mathcal{F}(x,x')+\beta V(x')\right\}
\]
Policy function:
\[
g(x)=\arg\max_{x'\in\Gamma(x)}\left\{\mathcal{F}(x,x')+\beta V(x')\right\}
\]

\begin{itemize}
  \item $V$: value of being in state $x$.
  \item $g$: optimal one-step-ahead decision rule.
\end{itemize}
\end{frame}

\begin{frame}{Standard NGM in Recursive Form (P4)}
\[
V(k)=\max_{k'\in\Gamma(k)}\left\{u\left(f(k)+(1-\delta)k-k'\right)+\beta V(k')\right\}
\]
with
\[
\Gamma(k)=[0,f(k)+(1-\delta)k]
\]

Current state: $k$. Control: $k'$ (equivalently consumption).
\end{frame}

\begin{frame}{How to Identify State Variables}
A variable should be a state if it is:
\begin{itemize}
  \item Predetermined at decision time.
  \item Payoff relevant for today and future.
  \item Not redundant given other states.
\end{itemize}

\begin{block}{Common mistake}
Leaving out a relevant predetermined variable breaks stationarity and yields wrong policy rules.
\end{block}
\end{frame}

\begin{frame}{Periodic Productivity Example (Chapter 4.4.2)}
If technology alternates deterministically $A_h, A_l$, recursive form uses two value functions:
\[
V_h(k)=\max_{k'\in\Gamma_h(k)}\{u(A_hf(k)+(1-\delta)k-k')+\beta V_l(k')\}
\]
\[
V_l(k)=\max_{k'\in\Gamma_l(k)}\{u(A_lf(k)+(1-\delta)k-k')+\beta V_h(k')\}
\]

State must include both $k$ and the regime indicator.
\end{frame}

\begin{frame}{Delayed Depreciation Example (Chapter 4.4.2)}
If capital equals two lagged investments, $k_t=i_{t-1}+i_{t-2}$, then two states are needed.

\medskip
Recursive representation uses $V(i',i)$ and choice $i''$.

\begin{block}{Intuition}
Non-standard technology often increases state dimension.
\end{block}
\end{frame}

\begin{frame}{Core Properties of the Value Function}
Under standard assumptions:
\begin{enumerate}
  \item Bellman operator is a contraction ($\beta<1$).
  \item $V$ is increasing when primitives are monotone.
  \item $V$ is concave when primitives/feasible set are concave-convex.
  \item Interior points allow differentiability of $V$.
  \item Policy function $g$ is increasing.
\end{enumerate}
\end{frame}

\begin{frame}{Value Iteration Algorithm}
\hypertarget{m-a8}{}
Start from any bounded continuous $V_0$ (often $V_0=0$), then iterate:
\[
V_{n+1}(x)=\max_{x'\in\Gamma(x)}\{\mathcal{F}(x,x')+\beta V_n(x')\}
\]
Error shrinks geometrically in sup norm:
\[
\|V_{n+1}-V\|\le\beta\|V_n-V\|
\]

\textbf{Computational message:} globally reliable convergence.
\apxbutton{a8}
\end{frame}

\begin{frame}{Example 4.7: Recursive Solution in Log-Cobb-Douglas Case}
With $u(c)=\log c$, $f(k)=Ak^\alpha$, $\delta=1$:
\[
V(k)=\max_{k'\ge 0}\{\log(Ak^\alpha-k')+\beta V(k')\}
\]

Value iteration reveals the policy converges to
\[
k'=\alpha\beta Ak^\alpha
\]

Same rule obtained from sequential approach.
\end{frame}

\begin{frame}{Functional Euler Equation}
\hypertarget{m-a9}{}
Bellman FOC (interior):
\[
-\mathcal{F}_2(x,x')=\beta V'(x')
\]
Envelope condition:
\[
V'(x)=\mathcal{F}_1(x,g(x))
\]
Combining:
\[
-\mathcal{F}_2(x,g(x))=\beta\mathcal{F}_1(g(x),g(g(x)))
\]
This is the functional Euler equation.
\apxbutton{a9}
\end{frame}

\begin{frame}{Why Functional Euler Is Useful but Not Enough}
\begin{itemize}
  \item It is a necessary condition for optimal policy.
  \item It can admit multiple candidate functions.
  \item Only one candidate attains the Bellman maximum globally.
\end{itemize}

\begin{block}{Practice}
Use functional Euler for intuition/local approximations, but verify with Bellman optimality.
\end{block}
\end{frame}

%% ============================================================
\section{4.4.6 Dynamics and Convergence}
%% ============================================================

\begin{frame}{Policy Dynamics in the Optimizing NGM}
Given $k_{t+1}=g(k_t)$, dynamics are a first-order system in the state.

\medskip
Steady states solve
\[
\bar{k}=g(\bar{k})
\]
Local approximation around steady state:
\[
g(k)\approx a_0+a_1(k-\bar{k})
\]

Convergence speed is governed by slope $a_1$.
\end{frame}

\begin{frame}{Stable vs Unstable Local Roots}
Chapter 4.4.6 shows a second-order condition can yield two local roots for slope.

\begin{itemize}
  \item One root satisfies $0<a_1<1$: stable, monotone convergence.
  \item Another root exceeds $1$: unstable, violates optimality/Bellman maximum.
\end{itemize}

\begin{block}{Economic Selection}
The Bellman problem chooses the stable policy branch.
\end{block}
\end{frame}

\begin{frame}{Calibration Insight from Local Slope}
Define curvature term near steady state (notation of chapter):
\[
\Theta\equiv\beta^2\frac{u' f''}{u''}>0
\]
Then local slope $a_1$ solves a quadratic expression.

\medskip
Interpretation:
\begin{itemize}
  \item $a_1\approx 0$: very fast convergence.
  \item $a_1\approx 1$: very slow convergence.
\end{itemize}

This replaces fixed-$s$ Solow speed with preference-dependent dynamics.
\end{frame}

\begin{frame}{Solow vs Optimizing NGM: Key Dynamic Difference}
\begin{columns}[T]
\column{0.48\textwidth}
\textbf{Solow}
\begin{itemize}
  \item $s$ exogenous constant.
  \item Dynamics from technology + depreciation + fixed saving rule.
\end{itemize}

\column{0.48\textwidth}
\textbf{Optimizing NGM}
\begin{itemize}
  \item $s_t=s(k_t)$ endogenous.
  \item Dynamics shaped by $u$, $f$, and expectations.
\end{itemize}
\end{columns}

\medskip
Same long-run logic, richer transition mechanism.
\end{frame}

%% ============================================================
\section{Wrap-up}
%% ============================================================

\begin{frame}{What Chapter 4 Delivers}
\begin{enumerate}
  \item A unified language for dynamic choices in macroeconomics.
  \item Euler equations as central intertemporal optimality conditions.
  \item Clear distinction between finite and infinite horizon boundary logic.
  \item nPg and TVC as fundamental infinite-horizon restrictions.
  \item Recursive methods that scale to quantitative and stochastic models.
\end{enumerate}
\end{frame}

\begin{frame}{Takeaways for Research Practice}
\begin{itemize}
  \item Always verify well-posedness before solving.
  \item Use sequential methods for theory and derivations.
  \item Use recursive methods for computation and policy functions.
  \item Check TVC/nPg consistency when proposing candidate paths.
  \item Keep growth consistency in mind: CRRA is structural, not cosmetic.
\end{itemize}
\end{frame}

\begin{frame}{Next Lecture (Chapter 5)}
We move from individual optimization to dynamic equilibrium:
\begin{itemize}
  \item Firms, households, and price determination.
  \item Decentralization of planner allocations.
  \item Competitive equilibrium as a dynamic system.
\end{itemize}
\end{frame}

%% ============================================================
\appendix
\section{Appendix Derivations}
%% ============================================================

\begin{frame}{A1. From Generic Constraint to Sequence-Only Problem}
\hypertarget{a1}{}
If $x_{t+1}=h(x_t,y_t)$ can be inverted to $y_t=\hat{h}(x_t,x_{t+1})$, define
\[
\mathcal{F}(x_t,x_{t+1})\equiv \hat{\mathcal{F}}\big(\hat{h}(x_t,x_{t+1})\big)
\]
Then (P1) becomes
\[
\max_{\{x_{t+1}\}}\sum_{t=0}^{T}\beta^t\mathcal{F}(x_t,x_{t+1})
\]
with $x_{t+1}\in\Gamma(x_t)$.

This is the standard form used for Euler equations in Chapter 4.
\backbutton{m-a1}
\end{frame}

\begin{frame}{A2. Two-Period Euler Equation with Kuhn-Tucker}
\hypertarget{a2}{}
FOCs:
\[
u'(c_0)-\mu-\lambda=0,
\qquad
\beta u'(c_1)-\frac{\mu}{1+r}=0
\]
If interior ($\lambda=0$):
\[
\mu=u'(c_0),\qquad \mu=\beta(1+r)u'(c_1)
\]
Hence
\[
u'(c_0)=\beta(1+r)u'(c_1)
\]
\backbutton{m-a2}
\end{frame}

\begin{frame}{A3. Finite-Horizon NGM Euler Equation}
\hypertarget{a3}{}
From
\[
\mathcal{L}=\sum_{t=0}^{T}\beta^t\left[u(c_t)+\mu_tk_{t+1}\right]
\]
where $c_t=f(k_t)+(1-\delta)k_t-k_{t+1}$, FOC for $k_{t+1}$ gives
\[
-u'(c_t)+\beta u'(c_{t+1})\left[f'(k_{t+1})+1-\delta\right]+\mu_t=0
\]
At interior dates, $\mu_t=0$.
\backbutton{m-a3}
\end{frame}

\begin{frame}{A4. EIS Under CRRA}
\hypertarget{a4}{}
Euler with gross return $R_{t,t+s}$:
\[
u'(c_t)=\beta^s u'(c_{t+s})R_{t,t+s}
\]
CRRA derivative: $u'(c)=c^{-\sigma}$, so
\[
\frac{c_{t+s}}{c_t}=\left(\beta^sR_{t,t+s}\right)^{1/\sigma}
\]
Take logs and differentiate w.r.t. $\log R_{t,t+s}$:
\[
EIS=\frac{1}{\sigma}
\]
\backbutton{m-a4}
\end{frame}

\begin{frame}{A5. Consolidating Budget Constraints with nPg}
\hypertarget{a5}{}
From $c_t+a_{t+1}=(1+r)a_t$, iterate forward:
\[
\sum_{t=0}^{T}\frac{c_t}{(1+r)^t}=a_0(1+r)-\frac{a_{T+1}}{(1+r)^T}
\]
Take limits and use
\[
\lim_{T\to\infty}\frac{a_{T+1}}{(1+r)^T}\ge 0
\]
obtain
\[
\sum_{t=0}^{\infty}\frac{c_t}{(1+r)^t}\le a_0(1+r)
\]
\backbutton{m-a5}
\end{frame}

\begin{frame}{A6. Proposition 4.4 Logic (Sketch)}
\hypertarget{a6}{}
Given concavity, for any feasible comparison path,
\[
\mathcal{F}(x_t,x_{t+1})-\mathcal{F}(x_t^*,x_{t+1}^*)
\]
is bounded above by tangent terms.

\medskip
Summing and applying Euler equations collapses interior terms.

\medskip
Remaining boundary term vanishes by TVC, implying candidate path gives weakly higher value than any feasible alternative.
\backbutton{m-a6}
\end{frame}

\begin{frame}{A7. Bellman Equation Derivation}
\hypertarget{a7}{}
Start from
\[
V(x_0)=\max_{\{x_{t+1}\}}\sum_{t=0}^{\infty}\beta^t\mathcal{F}(x_t,x_{t+1})
\]
Split first period:
\[
V(x_0)=\max_{x_1\in\Gamma(x_0)}\left\{\mathcal{F}(x_0,x_1)+\beta
\max_{\{x_{t+2}\}}\sum_{t=0}^{\infty}\beta^t\mathcal{F}(x_{t+1},x_{t+2})\right\}
\]
Bracketed term is $V(x_1)$ by stationarity.
\backbutton{m-a7}
\end{frame}

\begin{frame}{A8. Why Contraction Implies Reliable Iteration}
\hypertarget{a8}{}
Bellman operator $\mathcal{T}$ satisfies
\[
\|\mathcal{T}V-\mathcal{T}W\|\le\beta\|V-W\|,\qquad \beta<1
\]
Banach fixed-point theorem implies:
\begin{itemize}
  \item Unique fixed point $V$.
  \item Iteration $V_{n+1}=\mathcal{T}V_n$ converges to $V$ from any $V_0$.
  \item Geometric convergence rate $\beta$.
\end{itemize}
\backbutton{m-a8}
\end{frame}

\begin{frame}{A9. Functional Euler from Bellman + Envelope}
\hypertarget{a9}{}
Bellman FOC at interior optimum $x'=g(x)$:
\[
\mathcal{F}_2(x,g(x))+\beta V'(g(x))=0
\]
Envelope:
\[
V'(x)=\mathcal{F}_1(x,g(x))
\]
Substitute envelope one step ahead:
\[
-\mathcal{F}_2(x,g(x))=\beta\mathcal{F}_1(g(x),g(g(x)))
\]
\backbutton{m-a9}
\end{frame}

\begin{frame}{A10. Finite-Horizon Closed Form (Example 4.1)}
With log utility, Cobb-Douglas, and $\delta=1$, Euler gives recursion:
\[
\frac{1}{Ak_t^\alpha-k_{t+1}}=\beta\frac{\alpha A k_{t+1}^{\alpha-1}}{Ak_{t+1}^\alpha-k_{t+2}}
\]
Using terminal condition $k_{T+1}=0$ and backward induction yields
\[
k_{t+1}=\alpha\beta\frac{1-(\alpha\beta)^{T-t}}{1-(\alpha\beta)^{T-t+1}}Ak_t^{\alpha}
\]
\backbutton{m-a3}
\end{frame}

\begin{frame}{A11. Infinite-Horizon Log Consumption-Saving (Example 4.3)}
FOCs under consolidated budget:
\[
\beta^t\frac{1}{c_t}=\lambda\left(\frac{1}{1+r}\right)^t
\]
Thus
\[
c_t=[\beta(1+r)]^t c_0
\]
Plug into budget:
\[
\sum_{t=0}^{\infty}\beta^t c_0=a_0(1+r)
\]
Hence $c_0=a_0(1-\beta)(1+r)$.
\backbutton{m-a5}
\end{frame}

\begin{frame}{A12. Solow Link in Log-Cobb-Douglas NGM}
From Euler and guess $k'=sAk^\alpha$ with $\delta=1$:
\[
\frac{1}{Ak^\alpha-k'}=\beta\frac{\alpha A (k')^{\alpha-1}}{Ak'^\alpha-k''}
\]
Consistency of the guess implies
\[
s=\alpha\beta
\]
so policy is $k'=\alpha\beta Ak^\alpha$.
\backbutton{m-a6}
\end{frame}

\begin{frame}{A13. Steady State Condition in Stationary NGM}
At steady state, $k_{t+1}=k_t=\bar{k}$ and $c_{t+1}=c_t$.
Euler simplifies to
\[
1=\beta\left[f'(\bar{k})+1-\delta\right]
\]
Then resources imply
\[
\bar{c}=f(\bar{k})-\delta\bar{k}
\]
\backbutton{m-a6}
\end{frame}

\begin{frame}{A14. nPg vs TVC in One Line}
\begin{itemize}
  \item nPg: excludes paths with exploding debt in present value.
  \item TVC: excludes paths with leftover positive asset value at infinity.
\end{itemize}

\medskip
Feasibility + optimality jointly select economically meaningful paths in infinite horizon.
\backbutton{m-a6}
\end{frame}

\end{document}
